\subsection{Struktura promptu}
\label{subsec:prompts-structure}

Prompt je textový vstup, který definuje, co má jazykový model provést.
V kontextu analýzy studentského kódu v systému Kelvin se prompt skládá z několika logicky oddělených částí, jejichž správná struktura výrazně ovlivňuje kvalitu generovaného výstupu~\cite{white-prompt}.

\begin{enumerate}
    \item \textbf{Systémový prompt} obsahuje instrukce definující roli modelu, požadovaný jazyk odpovědi, formát výstupu a pravidla, která má model při analýze dodržovat.
    Typicky se jedná o pokyny jako \uv{Jsi zkušený programátor analyzující studentský kód v jazyce C++} doplněné o specifické požadavky na strukturu odpovědi.
    Systémový prompt je konstantní pro všechna odevzdání a jeho rozsah se pohybuje v řádu 50 až 200 řádků.

    \item \textbf{Zadání úlohy} poskytuje modelu kontext potřebný pro posouzení, zda studentské řešení odpovídá požadavkům.
    Bez zadání může model hodnotit pouze obecnou kvalitu kódu, ale není schopen posoudit, zda řešení splňuje konkrétní podmínky.

    \item \textbf{Zdrojový kód studenta} je hlavní analyzovaný vstup.
    Model obdrží kompletní obsah odevzdaného souboru včetně komentářů a prázdných řádků.
    Zachování těchto řádků je důležité, protože bez nich by model špatně odkazoval na konkrétní místa v kódu a výstup by byl matoucí.

    \item \textbf{Specifikace výstupního formátu} definuje, v jaké struktuře má model vrátit výsledky analýzy.
    Pro strojové zpracování na straně systému Kelvin je vhodný strukturovaný formát JSON, který obsahuje seznam nalezených problémů s uvedením čísla řádku, popisu problému a důležitosti.
    Tato specifikace je součástí systémového promptu a zajišťuje, že výstup modelu lze programově zpracovat bez dodatečného parsování volného textu.
\end{enumerate}

Správné oddělení těchto částí je důležité nejen pro kvalitu výstupu, ale i pro údržbu systému.
Pokud je systémový prompt oddělen od zadání a zdrojového kódu, lze jej snadno aktualizovat bez nutnosti měnit způsob, jakým systém Kelvin předává data modelu.
Příklad použitelného promptu je možné vidět v příloze~\ref{lst:simple-prompt}.

\endinput